\begin{frame}
\frametitle{Roteiro da Apresentação}
\tableofcontents
\end{frame}

\section{Simplex Padrão vs. Revisado}

\begin{frame}
\frametitle{Contextualização e Problema}
\begin{itemize}
    \item \textbf{O Problema:} Problemas reais (como escala de tripulação e redes de energia) possuem milhões de variáveis ($n$) e milhares de restrições ($m$).
    \item \textbf{Gargalo do Simplex Padrão:} 
    \begin{itemize}
        \item Atualiza toda a matriz $m \times n$ a cada iteração.
        \item Gera \textit{Fill-in}: Matrizes esparsas tornam-se densas, estourando a memória.
    \end{itemize}
    \item \textbf{Solução:} Método Simplex Revisado.
\end{itemize}
\end{frame}

\begin{frame}
\frametitle{Comparativo: Padrão vs. Revisado}
\begin{table}[]
\centering
\small
\begin{tabular}{p{0.45\textwidth}|p{0.45\textwidth}}
\hline
\textbf{Simplex Padrão (Tableau)} & \textbf{Simplex Revisado} \\ \hline
Calcula $B^{-1}a_j$ para \textbf{todas} as colunas não-básicas. & \textbf{Computação sob Demanda:} Calcula dados apenas da variável candidata. \\ \hline
Custo por iteração: $\approx O(m \cdot n)$. & Custo por iteração: $\approx O(m^2)$ (independe de $n$). \\ \hline
Erros numéricos se propagam por todo o quadro. & Erros restritos à inversa $B^{-1}$ (pode ser resetada). \\ \hline
\end{tabular}
\end{table}
\end{frame}

\section{O Método Simplex Revisado}

\begin{frame}
\frametitle{A Filosofia "Sob Demanda"}
No Revisado, trabalhamos com a matriz original $A$ e mantemos apenas a inversa da base $B^{-1}$.
\vspace{0.3cm}
\begin{block}{Passo 1: Multiplicadores Simplex (Dual)}
    Em vez de atualizar custos, calculamos $\pi$ resolvendo o sistema:
    \[ \pi^T = c_B^T B^{-1} \]
\end{block}
\begin{block}{Passo 2: Pricing (Quem entra?)}
    Calculamos o custo reduzido apenas para verificar a entrada:
    \[ \hat{c}_j = c_j - \pi^T a_j \]
    Se $\hat{c}_k > 0$, a variável $x_k$ entra na base.
\end{block}
\end{frame}

\begin{frame}
\frametitle{Algoritmo: Passos Finais}
\begin{block}{Passo 3: Coluna Atualizada}
    Calculamos apenas a coluna da variável vencedora ($a_k$):
    \[ \alpha = B^{-1} a_k \]
\end{block}

\begin{block}{Passo 4: Teste da Razão (Quem sai?)}
    Usamos a solução atual $\bar{b} = B^{-1}b$:
    \[ \theta = \min \left\{ \frac{\bar{b}_i}{\alpha_i} : \alpha_i > 0 \right\} \]
\end{block}

\begin{itemize}
    \item \textbf{Atualização:} A nova inversa $B^{-1}_{nova}$ é obtida multiplicando a anterior por uma matriz elementar (Eta).
\end{itemize}
\end{frame}